\section{The Boundaries}

Four things are regularly mistaken for trustful surrender: quietism, the pure-love doctrine condemned in Fénelon, fatalism, and presumption. The first two were condemned by name and in propositions, which makes them the most useful boundary markers available: instead of arguing about tendencies, one can read what was actually forbidden. The other two are errors of reasoning rather than condemned systems, and they are treated after the propositions.

\subsection{Molinos and \emph{Coelestis Pastor}}

Miguel de Molinos, a Spanish priest resident in Rome, published \emph{La Guía espiritual} in 1675; it was approved by theologians, went through some twenty editions in twelve years, and made its author one of the most sought-after directors in the city. In 1685 the Holy Office ordered his arrest. A decree of 28 August 1687 and the constitution \emph{Coelestis Pastor} condemned sixty-eight propositions that he had acknowledged as his own; the sentence was read publicly at Santa Maria sopra Minerva on 3 September 1687, and he abjured and was imprisoned for life.\endnote{Latin text of the condemned propositions and of the censure read 2026-07-25 in the web presentation of Denzinger's \emph{Enchiridion symbolorum} at \texttt{patristica.net}, which prints the revised (DS) numbering beside the older (DB) numbering; the Molinos propositions run DS 2201--2268 (DB 1221--1288), with the censure at DS 2269. An English translation of the constitution, including its narrative sections, was read the same day at \texttt{papalencyclicals.net}. \textbf{A discrepancy is recorded rather than resolved:} the Denzinger heading dates \emph{Coelestis Pastor} 20 November 1687, while the 1913 \emph{Catholic Encyclopedia} article on Molinos (\texttt{newadvent.org}, checked 2026-07-25) dates it 2 November 1687. No original printing of the constitution was consulted for this article; both witnesses are reported, and nothing in the argument depends on which date is right. Biographical details (the 1685 arrest, the 3 September 1687 reading of the sentence, the twenty editions of the \emph{Guía}) are as reported in that encyclopedia article. English renderings of the propositions in the table below are working glosses by this article; the Latin governs.}

What was condemned is not a mood but a chain of consequences. The first proposition sets the programme---\emph{oportet hominem suas potentias annihilare, et haec est via interna}, man must annihilate his powers, and this is the interior way---and the rest draw it out: to want to act is to offend God, who wills to be the sole agent, so one must abandon oneself wholly and afterwards remain \emph{velut corpus exanime}, like a lifeless body; natural activity is grace's enemy, ``because God wills to work in us without us''; the soul is not to seek its own salvation, nor to ask anything of God, nor to give thanks; temptations are met with no more than negative resistance; the virtues are to be abandoned; acts of love toward the Blessed Virgin, the saints, and the humanity of Christ are excluded as sensible; obscene external acts performed under diabolical ``violence'' are not sins and need not be confessed; through acquired contemplation one reaches a state of sinning no more; and the interior life is separated from confession, from directors' external judgement, and from the bishop's authority.\endnote{Propositions cited in this paragraph, in Denzinger's revised numbering: DS 2201 (1), 2202 (2), 2204 (4), 2212 (12), 2214 (14), 2215 (15), 2217 (17), 2231 (31), 2235 (35), 2241 (41), 2246--2248 (46--48), 2257 (57), 2259 (59), 2265--2267 (65--67). Latin read 2026-07-25 at the patristica.net witness cited above.}

The censure appended to the list is itself instructive, because it is graded rather than uniform: certain propositions are branded heretical (among them numbers 3, 13--15, and the whole block 41--53 on the ``violences''), others as proximate to heresy or savouring of heresy, others as erroneous, scandalous, blasphemous, offensive to pious ears, rash, or subversive of ecclesiastical jurisdiction, ``respectively''---that is, each qualification attaching to the propositions listed against it.\endnote{DS 2269, read 2026-07-25 at the patristica.net witness: ``Quas quidem propositiones tamquam haereticas (3 13-15 41-53), suspectas \ldots\ et erroneas \ldots\ scandalosas \ldots\ blasphemas \ldots\ piarum aurium offensivas \ldots\ temerarias \ldots\ christianae disciplinae relaxativas \ldots\ et eversivas (iurisdictionis eccl.: 68) et seditiosas (65) respective\ldots\ damnavimus.'' The grading matters: a blanket assertion that ``all sixty-eight were condemned as heretical'' misstates the act.} Reading the grading is part of reading the document: the condemnation is a juridical instrument with internal distinctions, not a thunderclap.

\subsection{Fénelon and \emph{Cum alias}}

Twelve years later the brief \emph{Cum alias ad apostolatus} of Innocent XII, dated 12 March 1699, condemned twenty-three propositions extracted from Fénelon's \emph{Explication des maximes des saints sur la vie intérieure}. The subject is different from Molinos's. Fénelon is not teaching the annihilation of the faculties or the suspension of the moral law; he is defending disinterested love, and the propositions that gave offence concern hope.\endnote{Latin text read 2026-07-25 at the patristica.net Denzinger witness: Innocent XII, brief \emph{Cum alias ad apostolatus}, 12 March 1699, \emph{Errores Francisci de Fénelon de amore erga Deum}, DS 2351--2374 (DB 1327--1349). English renderings below are working glosses by this article.}

The first proposition posits ``a habitual state of love of God which is pure charity and without any admixture of the motive of one's own interest,'' in which ``neither the fear of punishments nor the desire of rewards has any longer a part''; the second, that in the contemplative or unitive state ``every interested motive of fear and hope is lost''; the sixth, that in the state of holy indifference ``we no longer will salvation as our own salvation, as eternal deliverance, as the reward of our merits''; the eighth and tenth, that in the extreme trials the ordinarily conditional sacrifice of one's beatitude becomes in some manner absolute, the soul expiring with Christ on the cross in an ``involuntary impression of despair'' and there making ``the absolute sacrifice of its own interest with respect to eternity''; the ninth, that in such trials the soul can be invincibly persuaded, by a reflex persuasion, ``that it is justly reprobated by God''; the twenty-first, that ``the holy mystics excluded from the state of transformed souls the exercises of the virtues.''\endnote{DS 2351, 2352, 2356, 2358, 2359, 2360, 2371; Latin read 2026-07-25 at the witness cited above.}

The censure is as important as the list. The brief condemns the propositions as ``rash, scandalous, ill-sounding, offensive to pious ears, pernicious in practice, and even erroneous, respectively,'' and prohibits the book because its reading could lead the faithful gradually into errors already condemned. It does not call them heretical, and it does not call their author one.\endnote{DS 2374, read 2026-07-25: ``\ldots temerarias (1s 8 10 15-20 22), scandalosas (7 10 12 19-21), male sonantes (4-6 23), piarum aurium offensivas (8 18), in praxi perniciosas (2 14 17) ac etiam erroneas (1-7 10s 13 17-19 22s) respective, tenore praesentium damnamus et reprobamus.'' The absence of \emph{haereticas} from this list, and its presence in the Molinos censure, is a material difference between the two acts.} The distance between the two documents is thus considerable, and it is a distance that should be preserved by anyone who invokes them. Molinos was condemned for a system that dissolved the moral life; Fénelon for pressing one strand of an accepted tradition---the priority of charity over self-interest---to a point at which it began to exclude the theological virtue of hope. He submitted, read the condemnation from his own cathedral pulpit, and died in communion.

The Salesian sentence quoted earlier shows why the line is fine and why it is nonetheless a line. De Sales says the indifferent soul would run to hell if it knew God preferred that---but he says it under an explicit ``supposition of an impossible thing,'' a counterfactual whose antecedent he denies, and in a chapter whose whole burden is that God's will is the one thing loved. Fénelon's condemned propositions convert the impossible supposition into a state actually reached in extreme trials, in which the sacrifice becomes ``in some manner absolute'' and the soul acquiesces in a persuasion of its own reprobation. A rhetorical impossibility has become a spiritual programme; and a programme in which hope is suspended is a programme in which one of the three theological virtues has been asked to stand down.

\begin{fourstable}{Condemned proposition (gloss)}{Source}{What it denies}{How trustful surrender differs}
Man must annihilate his powers; this is the interior way. & Molinos, DS 2201 (prop.~1) & That the created faculties are goods to be perfected, and that God's action perfects rather than replaces them. & Grace perfects nature; God ``works in things in such a manner that things have their proper operation'' (ST I, q.~105, a.~5). Surrender submits the will's \emph{use}, not its existence.\\
To want to act is to offend God, who wills to be the sole agent; one must then remain like a lifeless body. & Molinos, DS 2202 (prop.~2) & Real secondary causality, and the dignity of acting that God confers. & ``If God governed alone, things would be deprived of the perfection of causality'' (ST I, q.~103, a.~6, ad~2); CCC 306 calls the use of creatures a token of God's greatness.\\
Natural activity is grace's enemy, for God wills to work in us without us. & Molinos, DS 2204 (prop.~4) & Co-operation with grace; the whole ascetical tradition of effort. & The soul consents, acts, and labours; abandonment follows ``after having done all that depends upon our efforts'' (Ramière, 1887 preface).\\
He who gives his will to God should care for nothing---not hell, not heaven, not his own salvation, whose hope he should remove. & Molinos, DS 2212 (prop.~12) & The theological virtue of hope and the duty of seeking salvation. & Surrender is an act of hope, not its suspension; the desire of beatitude is God's own design for the creature.\\
It is unfitting for the resigned to ask anything of God; petition is an imperfection. & Molinos, DS 2214--2215 (props.~14--15) & Petitionary prayer and thanksgiving, commanded by Christ. & Colombière's own compiled chapter is largely about petition: we obtain little because we ask too little.\\
Temptations require only negative resistance; if nature is aroused, let it be aroused. & Molinos, DS 2217 (prop.~17) & The obligation to resist temptation and to elicit contrary acts. & Acceptance concerns events already permitted and real; it never extends to consenting to sin, one's own or another's.\\
No meditative person exercises true interior virtues; the virtues must be abandoned. & Molinos, DS 2231 (prop.~31) & That charity is exercised in the acts of the virtues. & Patience itself is a virtue caused by charity and impossible without grace (ST II-II, q.~136, a.~3); abandonment is a mode of the virtues, not a substitute.\\
Interior souls should not elicit acts of love for the Blessed Virgin, the saints, or the humanity of Christ. & Molinos, DS 2235 (prop.~35) & The Incarnation's permanent place in Christian prayer. & Ramière's own criticism of the treatise he edited was that it did not stress enough that God's sanctifying action never occurs apart from the person of Jesus Christ.\\
External sinful acts done under diabolical ``violence'' are not sins; it is holier not to confess them. & Molinos, DS 2241, 2246--2248 (props.~41, 46--48) & Imputability, the moral law, and the sacrament of penance. & The gravest of the censures fell here (heretical). Surrender never suspends the moral law or the ordinary means of grace.\\
Through acquired contemplation one reaches the state of committing no further sin, mortal or venial. & Molinos, DS 2257 (prop.~57) & The universal need of vigilance and the Church's teaching on the impossibility of impeccability in this life. & Surrender presupposes that the surrendered man remains a sinner in need of the sacraments.\\
The soul in mystical death can will nothing but what God wills, because it no longer has a will: God has taken it away. & Molinos, DS 2261 (prop.~61) & The persistence of the created will in union. & Union is the conformity of a living will, not its removal; the ``death'' of self-will in Catholic usage is metaphorical.\\
Obedience binds only in the exterior; in the interior only God and the director enter, and the interior need not be manifested. & Molinos, DS 2265--2267 (props.~65--67) & Ecclesial accountability; the whole logic of an examinable spiritual life. & An abandonment that exempts itself from correction has removed the only external check on self-deception.\\
There is a habitual state of pure love without any admixture of the motive of self-interest, in which neither fear of punishment nor desire of reward has a part. & Fénelon, DS 2351 (prop.~1) & That hope remains a theological virtue in the perfect. & Charity is first; it does not expel hope. ``How could one of the theological virtues be contrary to another, and charity exclude hope?'' (Ramière).\\
In holy indifference we no longer will salvation as our own salvation or as the reward of our merits. & Fénelon, DS 2356 (prop.~6) & The rightly ordered desire of beatitude. & God wills to be glorified through the creature's happiness; the two interests are ordered, never separated.\\
In the extreme trials the sacrifice of one's beatitude becomes in some manner absolute; the soul may be persuaded that it is justly reprobated, and there makes an absolute sacrifice of its interest for eternity. & Fénelon, DS 2358--2360 (props.~8--10) & The permanence of hope in the darkest trial. & De Sales's parallel sentence is governed by an express ``supposition of an impossible thing''; a counterfactual is not a state to be entered.\\
The holy mystics excluded the exercises of the virtues from the state of transformed souls. & Fénelon, DS 2371 (prop.~21) & The continuity of the moral life in mystical states. & The saints of the tradition practise the virtues more, not less, as they advance.\\
\end{fourstable}

\subsection{Fatalism}

Fatalism and surrender agree on one sentence---that what happens was going to happen---and differ on everything the sentence means. Three differences are decisive.

The first concerns the terminus. Fate, in its ancient sense, is a necessitating series of causes, indifferent to the persons caught in it; providence is the plan of an intellect that loves, and its terminus is not an outcome but an end, the divine goodness communicated to creatures. That is why Aquinas can concede a defensible sense to the word ``fate''---the disposition of second causes as ordered by providence---and still record that the holy doctors avoided the term, and endorse Augustine's ruling that whoever means by fate the will of God may keep his opinion and hold his tongue.\endnote{ST I, q.~116, aa.~1--2, as cited in the preceding section.}

The second concerns contingency. Fatalism holds that everything happens necessarily; the doctrine of providence holds that some things happen necessarily and others contingently, and that both happen infallibly under providence \emph{in the mode proper to them}, because providence prepared necessary causes for the first and contingent causes for the second.\endnote{ST I, q.~22, a.~4, corpus and ad~1, as cited above.} A surrender that says ``it had to be this way'' is fatalist; a surrender that says ``this is what has in fact happened, and it did not escape God'' is not.

The third concerns the past and the future. Surrender to providence is, in its primary application, an act about what has already occurred or is presently occurring---the event as permitted and already real. Applied to the future, it is an act of hope about God's fidelity, not a prediction of the outcome and not a reason to stop deliberating. The moment ``God's will be done'' becomes a premise in practical reasoning about what to do next, it has been converted into fate.

\subsection{Presumption}

The opposite error is subtler because it looks like trust. Aquinas defines presumption as immoderate hope, and distinguishes two forms: relying on one's own power for what exceeds it, and relying on God's power in a disordered way.\endnote{ST II-II, q.~21, a.~1, corpus; the article treats presumption as a vice opposed to hope by excess, and locates the second kind as a hope of obtaining from God what it is not fitting for God to give in the way expected.} The devotional literature of surrender is unusually exposed to the second form, because the language of trust supplies a pious description of every refusal of means: the man who declines the treatment, the diocese that does not investigate, the parent who will not intervene.

The test is whether the trust is being used to license the omission of something one is otherwise bound to do. God's promises are made to those who use the means he provides; asking him to supply for a duty that one has deliberately abandoned is not confidence but an attempt to make him a party to one's negligence. De Sales's sickbed rule is the practical form of the distinction: while the outcome is unknown, the signified will governs, and the signified will says to use the remedies.

\subsection{Passivity, Prudence, and the Duties of State}

The most common practical failure of surrender is not heresy but sequence: it applies the act of acceptance before the moment when there is anything to accept. Three rules keep the sequence right, and all three come from the sources already quoted rather than from modern nervousness.

First, while the divine good-pleasure is undeclared, it is not a source of practical guidance at all, since it is known only by events. The rule of action in the meantime is the signified will---commandments, duties of state, counsels prudently applied---and it is a full rule, not a stopgap.\endnote{\emph{Treatise}, Book~IX, ch.~VI, as cited: ``as long as it is unknown to us, we must keep as close as possible to the will of God which is already declared or signified to us.''}

Second, effort and abandonment are not rivals but consecutive: ``we are to omit nothing which is requisite to bring the work which God has put into our hands to a happy issue, yet upon condition that, if the event be contrary, we should lovingly and peaceably embrace it.'' What is surrendered is the outcome, which was never in our power; what is not surrendered is the care, which was commanded.\endnote{\emph{Treatise}, Book~IX, ch.~VI, as cited, with the distinction between having charge of a work and being responsible for its event.}

Third, where the harm comes from another's malice, orthodox abandonment expressly authorises resistance. Ramière's fourth rule is unambiguous: if it is in our power to avert the consequences of malice and injustice, and if our true interest and the divine glory require measures to that end, ``let us do so without departing from the practice of the holy virtue of abandonment.'' De Sales says the same in his own register: conformity to God's permission of another's sin obliges us not to acquiesce but to aid, admonish, and assist, and only afterwards to turn the mind elsewhere.\endnote{Ramière, 1887 preface, rule II.4, and \emph{Treatise}, Book~IX, ch.~VIII, both as cited above.}

It follows, and should be said in plain words because the counsel has been misused this way, that acceptance of an event as permitted by God is not consent to the intention of the person who caused it, not a reason to remain in danger, not a bar to flight, medical care, reporting a crime, seeking justice, or protecting a dependent, and not a duty of silence. A person in danger should seek competent help; nothing in this article is advice for such a case.

\subsection{The Compressed Locution}

One more boundary belongs here, because it concerns the very book with which this article began. Seventeenth-century devotional prose speaks compressedly: it says that if we fall ill God is the cause of our illness, that wealth and poverty alike come from him, that his is the hand that deals the blow of death. Read against the architecture set out earlier, those sentences are defensible---God is the cause of every positive actuality, and natural corruption falls under what he can will incidentally in willing the order of nature. Read as modern flat prose, they say that God gave you cancer in the way an assailant gives you a wound.

Three controls keep the reading honest. The text's own exception is the first. Both parts of the compiled book state, in their opening paragraphs, that nothing happens in the universe or in our lives except by God's willing or allowing it, and both immediately except sin from that statement---which is the same exception de Sales makes in his own words: ``Nothing, except sin, is done without that will of God which is called absolute, or will of good-pleasure.'' The clause is not a decorative qualification; it is the doctrine. The Catechism's rule of reading is the second: attributing actions to God without naming secondary causes is a way of teaching trust, not a denial that the secondary causes exist. And the third is the distinction between the two evils: God wills no evil of fault ever, and wills natural defect or punishment only incidentally, under a good---from which it does not follow, and must not be inferred, that a particular person's affliction is a particular punishment.\endnote{The devotional formula and its stated exception are reported, not quoted, from the openings of Part~I, ch.~I and Part~II, ch.~IV of the modern English translation, read 2026-07-25 as described in the textual section; the quoted sentence in the body is de Sales, \emph{Treatise}, Book~IX, ch.~I, in Mackey's public-domain English. ST I, q.~19, a.~9 and CCC 304 and 311 supply the controls. The inference from affliction to punishment is not merely uncharitable but unwarranted: cf.\ John 9:3.}
